probe 选择发生在主码扫描前，选择倒排列表，而非数据库向量保留候选。对 IVF 质心 $\mu_i$，选择器排序计算得到的 FP32 值路由分数：
\begin{equation}
 s_i=\|\mu_i\|_2^2-2\langle q',\mu_i\rangle,
 \qquad (s_i,i)\ \text{ordered lexicographically}.
 \label{eq:5-probe-score}
\end{equation}
省略共同查询范数不改变质心顺序。与主码平方距离表不同，这些分数可能为负。输出为
\[
 n_{\rm sel}=\min(n_{\rm probe},n_{\rm list}),\qquad
 I=\operatorname{Top}^{(s_i,i)}_{n_{\rm sel}}\{0,\ldots,n_{\rm list}-1\},
\]
并附列表长度前缀 offsets。

OpenAI v69 使用 TF32 旋转／质心 GEMM 和 FP32 累加。probe 精确性针对其计算分数；FP32 \texttt{fma} 不会消除此前舍入（\cref{sec:7-precision}）。

\subsection{Radix 选择与确定性排序}

一个 block 进行四趟从最高有效位开始的 radix。对前缀限制直方图 $H_t$ 和从一开始的剩余排名 $r_t$（$r_0=n_{\rm sel}$），
\[
 b_t=\min\Bigl\{b:\sum_{a\leq b}H_t[a]\geq r_t\Bigr\},\qquad
 r_{t+1}=r_t-\sum_{a<b_t}H_t[a].
\]
选中的 bin 扩展前缀。四个 bytes 确定 cutoff／等键数量。

将带符号零规范化为正零，再按下式映射 bits $u$：
\begin{equation}
 F(u)=\begin{cases}
       \mathop{\sim}u,&\text{if the sign bit is set},\\
       u\mathbin{|}2^{31},&\text{otherwise}.
      \end{cases}
 \label{eq:5-probe-key}
\end{equation}
对非 NaN 分数，分类可得单调性：
\[
 \begin{aligned}
 D_1<D_2<0&\Rightarrow u_1>u_2\Rightarrow{\sim}u_1<{\sim}u_2,\\
 0\leq D_1<D_2&\Rightarrow u_1<u_2\Rightarrow
 (u_1\mathbin{|}2^{31})<(u_2\mathbin{|}2^{31}),\\
 D_1<0\leq D_2&\Rightarrow F(u_1)<2^{31}\leq F(u_2).
 \end{aligned}
\]
该映射覆盖无穷／零，不同于仅适用于非负值的 \cref{eq:5-packed-key}。重复使用 \texttt{fma} 固定重算键。

直方图存储为
\[
 256\ \text{bins}\times32\ \text{lane replicas}\times4\ \text{bytes}=32\,\mathrm{KiB}.
\]
复制减少同 warp 冲突，但仍有 warp 间竞争／归约成本。

令 $h$ 为 cutoff 键。收集阶段形成
\[
 A_< =\{i:F(u_i)<h\},\quad A_= =\{i:F(u_i)=h\},\quad
 m=n_{\rm sel}-|A_<|.
\]
取 $A_<$ 及 $A_=$ 中 ID 最小的 $m$ 项，使用有序 ballot／前缀计数，而非原子到达顺序。对所得集合 $A$，
\[
 \operatorname{rank}(i)=|\{j\in A:(s_j,j)\prec(s_i,i)\}|,
 \qquad\mathrm{out}[\operatorname{rank}(i)]=i.
\]

选定长度产生前缀 offsets。对超出数量的 probes，
\[
 i_j=i_{n_{\rm sel}-1},\quad L_j=0\quad(j\geq n_{\rm sel}),
 \qquad\omega_{j+1}=\omega_j+L_j.
\]
零长度避免重复扫描。要求 $n_{\rm list},n_{\rm probe}>0$。

\subsection{成本与独立正确性约束}

直方图／收集重新读取分数；排名成本为 $O(n_{\rm sel}^2)$，共享内存随之增长。小规模选择也需支付固定直方图存储；这是确定性集成，不是统一的线性成本。

测试覆盖并列、零、无穷、极值和超量 probes（\cref{sec:7-4}）。

FP16 近似和跨进程检查另行处理（\cref{sec:7-4}）。
